\chapter{Bipolar Disorder}
\index{Bipolar Disorder}

\section*{Definition and Overview}
Bipolar disorder is a long-term mental health condition characterised by marked swings in mood, energy, and activity. People experience episodes of elevated or irritable mood, known as mania or hypomania, alternating with episodes of depression. Between episodes, many people function well, and the severity and frequency of episodes vary widely between individuals.

Bipolar disorder was formerly called manic depression. It is a chronic condition, but with mood-stabilising medicines, psychological support, and predictable routines, most people manage the illness effectively and live stable, productive lives.

\section*{Epidemiology}
Bipolar disorder affects roughly 1 to 2 per cent of the population worldwide, with men and women affected about equally.

\begin{itemize}
    \item \textbf{Age at onset:} the disorder usually begins in late adolescence or early adulthood, commonly between the ages of 15 and 25
    \item \textbf{Family history:} the condition is strongly familial, and people with a first-degree relative who has bipolar disorder are at markedly higher risk
    \item \textbf{Coexisting conditions:} anxiety disorders, attention deficit hyperactivity disorder, and substance use disorders are common
    \item \textbf{Diagnostic delay:} many people are not diagnosed for several years after their first mood episode
    \item \textbf{Universal occurrence:} bipolar disorder occurs in all cultural and socioeconomic groups
\end{itemize}

\section*{Aetiology and Pathophysiology}
The exact cause is not fully understood, but bipolar disorder is thought to arise from a combination of genetic and environmental factors.

\begin{itemize}
    \item \textbf{Genetics:} the disorder runs in families, and many different genes are believed to contribute to risk
    \item \textbf{Brain chemistry:} imbalances in neurotransmitters such as serotonin, dopamine, and noradrenaline, which regulate mood, energy, and motivation
    \item \textbf{Brain structure and function:} differences in the size and activity of brain regions involved in emotion regulation, reward, and impulse control
    \item \textbf{Stressful life events:} trauma, major loss, relationship breakdown, or significant sleep disruption can trigger episodes in susceptible people
    \item \textbf{Substance use:} alcohol and recreational drugs, particularly stimulants, can precipitate or worsen mood episodes
    \item \textbf{Circadian rhythm disruption:} disturbed sleep-wake patterns play an important role in triggering manic episodes
\end{itemize}

\section*{Clinical Features}
Symptoms occur in episodes and depend on whether the mood is elevated or depressed.

\begin{itemize}
    \item \textbf{Manic features:} elevated or irritable mood, increased energy, a reduced need for sleep, rapid and pressured speech, racing thoughts, inflated self-esteem, poor judgement, and risky behaviour such as excessive spending
    \item \textbf{Hypomanic features:} the same symptoms as mania but less severe, without psychosis or major functional impairment, and often experienced as a pleasant or productive state
    \item \textbf{Depressive features:} persistent sadness, loss of interest in usual activities, low energy, sleep and appetite changes, feelings of worthlessness, and poor concentration
    \item \textbf{Mixed episodes:} features of mania and depression occurring together, such as agitation combined with hopelessness
    \item \textbf{Psychotic features:} in severe episodes, false beliefs (delusions) or hallucinations may occur, usually consistent with the mood
    \item \textbf{Loss of insight:} during manic episodes people may not recognise that their behaviour is unusual or harmful
    \item \textbf{Functional impact:} episodes can disrupt work, study, finances, and relationships
\end{itemize}

\section*{Differential Diagnosis}
A careful assessment is needed because several conditions can resemble bipolar disorder.

\begin{itemize}
    \item \textbf{Major depressive disorder:} suspected when only depressive episodes have occurred and there is no history of mania or hypomania
    \item \textbf{Cyclothymic disorder:} milder and more chronic mood swings that do not meet the full criteria for mania or depression
    \item \textbf{Schizoaffective disorder and schizophrenia:} when psychotic symptoms are prominent
    \item \textbf{Anxiety disorders:} which may cause agitation, restlessness, and poor sleep
    \item \textbf{Attention deficit hyperactivity disorder:} overlapping features of impulsivity, distractibility, and restlessness
    \item \textbf{Substance-induced mood disturbance:} from alcohol, stimulants, or medicines such as corticosteroids
    \item \textbf{Medical conditions:} thyroid disease, neurological disorders, and other physical illnesses that affect mood
    \item \textbf{Borderline personality disorder:} rapid mood shifts that are often triggered by interpersonal events
\end{itemize}

\section*{When to Seek Medical Care}
Seek help from a general practitioner or mental health professional if mood swings are causing distress, affecting sleep or judgement, or interfering with work, study, or relationships. Anyone with new or worsening symptoms, or who is concerned that a relative is developing mania or depression, should be assessed early.

\textbf{Contact your local emergency services for:}
\begin{itemize}
    \item Thoughts of suicide, self-harm, or harming others
    \item Someone expressing a plan to end their life
    \item Severe agitation, confusion, or behaviour that is dangerous to the person or others
    \item Marked psychosis, or an inability to care for oneself safely
\end{itemize}

If you are feeling distressed, contact a local crisis service or your local mental-health crisis service.

\section*{Diagnosis and Investigations}
There is no single test for bipolar disorder; the diagnosis is based on careful clinical assessment.

\begin{itemize}
    \item \textbf{Clinical interview:} a detailed history of mood episodes, their timing, duration, and impact, ideally including information from family members who have observed the episodes
    \item \textbf{Mental state examination:} assessment of mood, speech, thought, perception, and cognition at the time of review
    \item \textbf{Mood charting:} keeping a diary of mood, sleep, energy, and activity over time can clarify the pattern
    \item \textbf{Questionnaires:} standardised mood-rating scales may support the assessment
    \item \textbf{Physical examination and blood tests:} thyroid function, vitamin B12, calcium, and other tests to exclude medical causes of mood change
    \item \textbf{Drug screen:} to identify substance-related mood disturbance
\end{itemize}

\section*{Management}
Bipolar disorder is managed over the long term, usually through a combination of approaches.

\begin{itemize}
    \item \textbf{Mood-stabilising medicines:} lithium and certain anticonvulsants are the mainstay of treatment, taken continuously to prevent relapses
    \item \textbf{Antipsychotic medicines:} used for acute mania or psychosis, and sometimes as long-term mood stabilisers
    \item \textbf{Antidepressants:} used for depressive episodes, usually together with a mood stabiliser and under close supervision, because they can trigger mania in some people
    \item \textbf{Psychological therapies:} psychoeducation, cognitive behaviour therapy, and family-focused therapy help with coping, adherence, and relapse prevention
    \item \textbf{Lifestyle regulation:} regular sleep, predictable routines, stress management, and avoidance of alcohol and recreational drugs
    \item \textbf{Crisis and hospital care:} required for severe mania, psychosis, or significant suicide risk
    \item \textbf{Allied health support:} occupational therapy, peer support groups, and case management help with social and occupational functioning
\end{itemize}

\section*{Complications}
\begin{itemize}
    \item \textbf{Suicide and self-harm:} the most serious complication of the illness
    \item Financial and relationship damage resulting from behaviour during manic episodes
    \item Alcohol and drug misuse, often used to cope with symptoms
    \item Occupational and educational disruption, including unemployment
    \item Side effects of treatment, including weight gain, tremor, thyroid dysfunction, and kidney effects with long-term lithium use
    \item Coexisting medical conditions, including anxiety, cardiovascular disease, and diabetes
\end{itemize}

\section*{Prognosis and Outcomes}
Bipolar disorder is a lifelong condition for most people and typically runs a relapsing and remitting course. With regular mood-stabilising treatment and good support, many people have long periods of stability and maintain full, productive lives.

Relapses tend to become less frequent and less severe in many cases, although poor adherence to medication is a common reason for recurrence. Early diagnosis, a consistent treatment team, and strong family and social support improve the outlook. Some people continue to experience significant functional difficulties even between episodes, so ongoing monitoring and support are important.

\section*{Prevention and Self-Care}
\begin{itemize}
    \item Take mood-stabilising medicines exactly as prescribed, and discuss any concerns with the doctor before stopping
    \item Maintain a regular sleep schedule and avoid shift changes where possible
    \item Keep a routine for meals, exercise, and daily activities
    \item Learn to recognise early warning signs of a mood episode, such as a reduced need for sleep
    \item Avoid alcohol and recreational drugs
    \item Reduce stress, and seek support early during difficult times
    \item Keep regular appointments with the general practitioner, psychiatrist, and treating team
    \item Involve trusted family or friends in recognising early changes in mood
\end{itemize}

\section*{Sources and Further Reading}
The following organisations provide reliable information on bipolar disorder for patients, families, and health professionals.

\begin{itemize}
    \item NHS. \textit{Information on bipolar disorder, its symptoms, and treatment.} https://www.nhs.uk/
    \item Mayo Clinic. \textit{Overview of bipolar disorder and its management.} https://www.mayoclinic.org/
    \item NICE. \textit{Clinical guidance on the assessment and management of bipolar disorder.} https://www.nice.org.uk/
    \item National Institute of Mental Health. \textit{Educational resources on bipolar disorder.} https://www.nimh.nih.gov/
    \item World Health Organization. \textit{Resources on mental health conditions, including bipolar disorder.} https://www.who.int/
\end{itemize}
